Frente a las plataformas anteriores, el SCDEE adopta varias decisiones de
diseño diferentes. Primero, las zonas de identificación del estudiante son
configurables por el creador del examen y admiten múltiples campos asociados
al examen. Segundo, la exportación de calificaciones se ofrece en un formato
flexible configurable por la organización y adecuado para su importación en
SIGMA Academic \cite{sigmaaie}, el sistema de gestión académica que utilizan
la propia \acs{uam} \cite{uamsigma} y otras universidades españolas.
Tercero, el sistema admite despliegue dual: cada organización elige entre
una instalación \textit{on-premise} y una operación en modo \acs{saas}
gestionada por un proveedor; la arquitectura de contenedores se mantiene
idéntica en ambos modos. Cuarto, la seguridad se integra en el diseño desde
el origen y no como una capa añadida: registro de auditoría inmutable de toda
operación sensible y protección anti-fuga de datos mediante marca de agua
dinámica con identidad del alumno, \acs{url} temporales con caducidad
corta y cifrado del \acs{pdf} en tránsito.

% La \ref{table:sota-platform-comparison} sintetiza la comparación.


% \begin{table}[Comparativa de plataformas y SCDEE]%
%             {table:sota-platform-comparison}%
%             {Comparativa de las plataformas de corrección digital
%              analizadas con el SCDEE.}
%   \small
%   \renewcommand{\arraystretch}{1.35}
%   \begin{tabular}{p{3.0cm} p{2.7cm} p{2.7cm} p{2.0cm} p{3.0cm}}
%     \hline
%     \textbf{Característica} &
%       \textbf{Crowdmark} &
%       \textbf{Gradescope} &
%       \textbf{Akindi} &
%       \textbf{SCDEE} \\
%     \hline
%     Tipo de respuesta evaluada &
%       Abierta, código, opción múltiple &
%       Abierta, código, opción múltiple &
%       Solo opción múltiple &
%       Abierta \\
%     Identificación del estudiante &
%       \acs{ocr} sobre zona fija de portada (ID numérico) y \acs{qr} &
%       \acs{ocr} y \textit{matching} automático contra lista &
%       \acs{omr} sobre código de estudiante &
%       \acs{qr} y \acs{ocr} de zonas configurables, con \textit{matching}
%       tolerante a errores \\
%     Asistencia al corrector con IA &
%       Limitada y opcional &
%       Agrupamiento de respuestas similares &
%       No &
%       No \\
%     Modelo de despliegue &
%       \acs{saas} &
%       \acs{saas} (regiones EU/CA/AU) &
%       \acs{saas} &
%       \acs{saas}, con \textit{on-premise} también soportado \\
%     Integración \acs{lms} nativa &
%       Canvas, Blackboard, Brightspace, Moodle, \acs{lti} 1.3 &
%       Canvas, Blackboard, Moodle, D2L, Sakai (\acs{lti} 1.3) &
%       Canvas, Blackboard &
%       No nativa en v1.0; trabajo futuro \\
%     \acs{sso} institucional &
%       Sí &
%       Sí (\acs{saml}) &
%       Vía \acs{lms} &
%       \Dfn{plugin} diseñado, no implementado en v1.0 \\
%     Integración con sistemas académicos institucionales &
%       No documentada &
%       No documentada &
%       No documentada &
%       Exportación de calificaciones en formato SIGMA \\
%     Origen &
%       Empresa privada (Toronto, 2014) &
%       \textit{UC Berkeley} (2014); Turnitin (2018) &
%       Empresa privada &
%       \acs{tfg} \acs{eps}-\acs{uam} (2026) \\
%     \hline
%   \end{tabular}
% \end{table}
